\section{第13章の解答}\label{解答:第13章}

\paragraph{13.1}
分母の次数より分子の次数が高いので,先に多項式の割り算をする.
\[
 x^3+x^2+2x+1=(x+1)(1+x^2)+x
\]
だから
\[
 \begin{aligned}
 \int_0^1\frac{x^3+x^2+2x+1}{1+x^2}\dd x
 &=\int_0^1\left(x+1+\frac{x}{1+x^2}\right)\dd x\\
 &=\left[\frac{x^2}{2}+x+\frac12\log(1+x^2)\right]_0^1\\
 &=\boxed{\frac32+\frac12\log2}.
 \end{aligned}
\]
余りの分子$x$が分母の微分$2x$の定数倍になったため,対数の形で積分できた.

\paragraph{13.2}
$u=1+x^2$と置くと$\dd u=2x\dd x$であり,\,$x=0,1$はそれぞれ$u=1,2$に対応する.従って
\[
 \int_0^1x\sqrt{1+x^2}\dd x
 =\frac12\int_1^2u^{\frac12}\dd u
 =\left[\frac13u^{\frac32}\right]_1^2
 =\boxed{\frac{2\sqrt2-1}{3}}.
\]
積分変数を$u$へ変えた時点で,端点も$u$の値へ変えている.

\paragraph{13.3}
$\log x$を微分し,\,$x$を積分すると,残りは整式の積分になる.部分積分により
\[
 \begin{aligned}
 \int_1^{\e}x\log x\dd x
 &=\left[\frac{x^2}{2}\log x\right]_1^{\e}
   -\frac12\int_1^{\e}x\dd x\\
 &=\frac{\e^2}{2}-\frac{\e^2-1}{4}
 =\boxed{\frac{\e^2+1}{4}}.
 \end{aligned}
\]

\paragraph{13.4}
$G(u)=\int_0^u\sqrt{1+t^2}\dd t$と置く.基本定理より$G'(u)=\sqrt{1+u^2}$であり,\,$F(x)=G(x^2)-G(x)$だから,合成関数の微分によって
\[
 F'(x)=2x\sqrt{1+x^4}-\sqrt{1+x^2}.
\]
上端だけでなく,下端からも項が現れる.\,$x=1$では上下端が一致するため$F(1)=0$,また$F'(1)=\sqrt2$である.従って接線は
\[
 \boxed{y=\sqrt2(x-1)}.
\]

\paragraph{13.5}
求める積分を$I$とする.\,$x=-t$と置き,\,$1+x^2$が偶関数であることを使えば
\[
 I=\int_{-a}^{a}\frac{1+t^2}{1+\e^{-t}}\dd t.
\]
元の式と足すと,
\[
 \frac1{1+\e^x}+\frac1{1+\e^{-x}}=1
\]
によって重みが消える.従って
\[
 2I=\int_{-a}^{a}(1+x^2)\dd x
 =2a+\frac{2a^3}{3},\qquad
 \boxed{I=a+\frac{a^3}{3}}.
\]
不定積分を求めずに値が決まるのは,区間と分子が反転に対して変わらないためである.

\paragraph{13.6}
交点は$x^2=x$より$x=0,1$である.\,$0\leqq x\leqq1$では$x^2\leqq x$なので,面積$S$は
\[
 S=\int_0^1(x-x^2)\dd x
 =\left[\frac{x^2}{2}-\frac{x^3}{3}\right]_0^1
 =\boxed{\frac16}.
\]
回転すると,\,$x$を固定した断面は外半径$x$,内半径$x^2$の円環になる.従って体積$V$は
\[
 V=\pi\int_0^1\{x^2-(x^2)^2\}\dd x
 =\pi\left(\frac13-\frac15\right)
 =\boxed{\frac{2\pi}{15}}.
\]
面積では上下の高さの差を積分し,体積では二つの半径の平方の差を積分している.

\paragraph{13.7}
積分の両端を微分すると
\[
 \begin{aligned}
 H'(a)&=\{(a+1)^2-2(a+1)\}^2-(a^2-2a)^2\\
 &=(a^2-1)^2-(a^2-2a)^2\\
 &=(2a-1)(2a^2-2a-1).
 \end{aligned}
\]
従って極値の候補は$a=\frac12,\frac{1-\sqrt3}{2},\frac{1+\sqrt3}{2}$である.候補の値を比べ,実数全体での最小値を確かめるため,積分区間の中心を使って式を整理する.

$y=a-\frac12$とし,\,$x=a+\frac12+s$と置くと,\,$-\frac12\leqq s\leqq\frac12$,\,$x-1=y+s$である.\,$x^2-2x=(x-1)^2-1$なので
\[
 H(a)=\int_{-\frac12}^{\frac12}\{(y+s)^2-1\}^2\dd s.
\]
展開後の$s$の奇数次の項は積分すると消える.残る項は
\[
 \begin{aligned}
 H(a)
 &=\int_{-\frac12}^{\frac12}
   \{s^4+(6y^2-2)s^2+(y^2-1)^2\}\dd s\\
 &=\frac1{80}+\frac{6y^2-2}{12}+(y^2-1)^2\\
 &=y^4-\frac32y^2+\frac{203}{240}
 =\left(y^2-\frac34\right)^2+\frac{17}{60}.
 \end{aligned}
\]
最後の平方は$0$以上で,\,$y=\pm\frac{\sqrt3}{2}$のときに限り$0$になる.従って
\[
 \boxed{\min H(a)=\frac{17}{60},\qquad
 a=\frac{1-\sqrt3}{2},\ \frac{1+\sqrt3}{2}}.
\]
二つの区間が同じ最小値を与えるのは,被積分関数が$x=1$に関して対称であることにも対応する.

\paragraph{13.8}
対数の中を$u=1+x^2$と置くと,\,$\dd u=2x\dd x$,\,$x^2=u-1$だから
\[
 x^3\dd x=\frac12(u-1)\dd u.
\]
区間は$[1,2]$となり,求める積分は$\frac12\int_1^2(u-1)\log u\dd u$である.ここで$\log u$を微分する部分積分を行う.
\[
 \begin{aligned}
 \frac12\int_1^2(u-1)\log u\dd u
 &=\frac12\left[\left(\frac{u^2}{2}-u\right)\log u\right]_1^2
   -\frac12\int_1^2\left(\frac u2-1\right)\dd u\\
 &=\left[-\frac{u^2}{8}+\frac u2\right]_1^2
 =\frac12-\frac38
 =\boxed{\frac18}.
 \end{aligned}
\]
端点の項は,\,$u=1$では$\log1=0$,\,$u=2$では$\frac{u^2}{2}-u=0$となるため,両方とも消える.

\paragraph{13.9}
一項を
\[
 \frac{k^2}{n(n^2+k^2)}
 =\frac1n\cfrac{\left(\cfrac kn\right)^2}{1+\left(\cfrac kn\right)^2}
\]
と書く.\,$f(x)=\frac{x^2}{1+x^2}$は$[0,1]$で連続なので,この和の極限はその区間での積分である.分子を$x^2=(1+x^2)-1$と分けると
\[
 \begin{aligned}
 \int_0^1\frac{x^2}{1+x^2}\dd x
 &=\int_0^1\left(1-\frac1{1+x^2}\right)\dd x\\
 &=\left[x-\arctan x\right]_0^1
 =\boxed{1-\frac{\pi}{4}}.
 \end{aligned}
\]

\paragraph{13.10}
各$x\in[0,1]$で,縦の線分は$y=x-1$から$y=x$まであり,常に$x$軸を含む.従って回転断面の中心には穴がなく,円板になる.端点の軸からの距離はそれぞれ$1-x$,\,$x$なので,半径$r(x)$は
\[
 r(x)=
 \begin{cases}
  1-x&\left(0\leqq x\leqq\frac12\right),\\
  x&\left(\frac12\leqq x\leqq1\right).
 \end{cases}
\]
切り替わりは$x=1-x$,すなわち$x=\frac12$で起こる.従って
\[
 \begin{aligned}
 V&=\pi\int_0^{\frac12}(1-x)^2\dd x
   +\pi\int_{\frac12}^{1}x^2\dd x\\
 &=2\pi\int_{\frac12}^{1}x^2\dd x
 =\frac{2\pi}{3}\left(1-\frac18\right)
 =\boxed{\frac{7\pi}{12}}.
 \end{aligned}
\]
上下の部分を別々に回して体積を足すと,同じ円板内の点を重複して数える.ここでは回転後にできる断面そのものを一枚ずつ数えた.

\paragraph{13.11}
まず
\[
 K_0=\int_0^1\frac1{1+x}\dd x=\log2.
\]
$n\geqq1$では,隣り合う二つの積分を足すと分母が消える.
\[
 K_n+K_{n-1}
 =\int_0^1\frac{x^n+x^{n-1}}{1+x}\dd x
 =\int_0^1x^{n-1}\dd x=\frac1n.
\]
従って漸化式は$K_n=\frac1n-K_{n-1}$であり,順に
\[
 K_1=1-\log2,\qquad
 K_2=\log2-\frac12,\qquad
 \boxed{K_3=\frac56-\log2}
\]
となる.

また$0\leqq x\leqq1$では$0\leqq\frac{x^n}{1+x}\leqq x^n$だから
\[
 0\leqq K_n\leqq\frac1{n+1}\longrightarrow0.
\]
ここで漸化式の添字を$k$とし,両辺に$(-1)^{k-1}$を掛けて$k=1,\ldots,n$について足す.
中間の$K_k$が消え,
\[
 \sum_{k=1}^n\frac{(-1)^{k-1}}k
 =K_0+(-1)^{n-1}K_n
 =\log2+(-1)^{n-1}K_n
\]
となる.右辺の最後の項は$0$へ近づくので,求める極限は$\log2$である.同時に,部分和と$\log2$との差の絶対値が$\frac1{n+1}$以下であることも分かる.

\paragraph{13.12}
積分を$x=1$で分け,\,$0$と無限大を独立に調べる.\,$x\to0+$では
\[
 \cfrac{\cfrac{x^\alpha}{1+x^2}}{x^\alpha}
 =\frac1{1+x^2}\longrightarrow1.
\]
従って$0$付近では$x^\alpha$との極限比較が使え,収束条件は$\alpha>-1$である.
一方,\,$x\to\infty$では
\[
 \cfrac{\cfrac{x^\alpha}{1+x^2}}{x^{\alpha-2}}
 =\frac{x^2}{1+x^2}\longrightarrow1.
\]
無限大で$x^{\alpha-2}$が積分可能になる条件は$\alpha-2<-1$,すなわち$\alpha<1$である.
従って全体の収束条件は
\[
 \boxed{-1<\alpha<1}.
\]
$\alpha=0$では,端点を避けて計算してから極限を取り,
\[
 \int_0^{\infty}\frac1{1+x^2}\dd x
 =\lim_{R\to\infty}\arctan R-\lim_{\varepsilon\to0+}\arctan\varepsilon
 =\boxed{\frac{\pi}{2}}.
\]

\paragraph{13.13}
最初から$f$が微分可能だと仮定せず,右辺の積分から確かめる.
\[
 A(x)=\int_0^xf(t)\dd t,\qquad
 B(x)=\int_0^xtf(t)\dd t
\]
と置くと,\,$f$が連続であるため$A'(x)=f(x)$,\,$B'(x)=xf(x)$である.与式は
\[
 f(x)=1+xA(x)-B(x)
\]
と書けるので,右辺が微分可能であり,
\[
 f'(x)=A(x)+xf(x)-xf(x)=A(x)
\]
を得る.さらに$f''(x)=f(x)$である.また$x=0$を代入すると
\[
 f(0)=1,\qquad f'(0)=0.
\]

ここからは,\,$f''=f$を満たす関数を積の微分によって求める.
\[
 \frac{\dd}{\dd x}\{\e^{-x}(f(x)+f'(x))\}
 =\e^{-x}(f''(x)-f(x))=0
\]
だから,括弧を含む積は定数である.\,$x=0$での値は$1$なので
$f'(x)+f(x)=\e^x$となる.今度は$\e^xf(x)$を微分すれば
\[
 \frac{\dd}{\dd x}\{\e^xf(x)\}=\e^x(f'(x)+f(x))=\e^{2x}.
\]
$0$から$x$まで積分し,\,$f(0)=1$を使うと
\[
 \e^xf(x)-1=\frac{\e^{2x}-1}{2},\qquad
 \boxed{f(x)=\frac{\e^x+\e^{-x}}2}.
\]

微分して得た条件だけでは,元の積分の等式を満たすかの確認が残る.
求めた関数を$g(x)=\frac{\e^x+\e^{-x}}2$と書くと,\,$g''=g$,\,$g(0)=1$,\,$g'(0)=0$である.従って部分積分によって
\[
 \begin{aligned}
 \int_0^x(x-t)g(t)\dd t
 &=\int_0^x(x-t)g''(t)\dd t\\
 &=\bigl[(x-t)g'(t)\bigr]_0^x+\int_0^xg'(t)\dd t\\
 &=-xg'(0)+g(x)-g(0)=g(x)-1.
 \end{aligned}
\]
よって元の等式も成立する.また,条件を満たすどの$f$も上の計算から同じ式になるため,解は一つに定まる.

\paragraph{13.14}
二つの制約を満たす一次関数を先に探す.\,$\ell(x)=a+bx$を代入すると
\[
 a+\frac b2=1,\qquad \frac a2+\frac b3=0.
\]
これを解くと$a=4$,\,$b=-6$なので,\,$\ell(x)=4-6x$である.
任意の条件を満たす$f$について$g(x)=f(x)-\ell(x)$と置けば
\[
 \int_0^1g(x)\dd x=0,\qquad
 \int_0^1xg(x)\dd x=0.
\]
従って
\[
 \int_0^1\ell(x)g(x)\dd x
 =4\int_0^1g(x)\dd x-6\int_0^1xg(x)\dd x=0.
\]
この関係を平方の展開に使うと
\[
 \begin{aligned}
 \int_0^1f(x)^2\dd x
 &=\int_0^1\ell(x)^2\dd x
   +2\int_0^1\ell(x)g(x)\dd x+\int_0^1g(x)^2\dd x\\
 &=\int_0^1(16-48x+36x^2)\dd x+\int_0^1g(x)^2\dd x\\
 &=4+\int_0^1g(x)^2\dd x\geqq4.
 \end{aligned}
\]
$g$は連続なので,最後の積分が$0$になるのは$g(x)=0$が全区間で成り立つ場合に限る.
実際,ある点で$g\ne0$なら,その近くでは$g^2$が正の一定値以上になり,積分も正になる.従って
\[
 \boxed{\min\int_0^1f(x)^2\dd x=4,\qquad f(x)=4-6x}.
\]
制約から候補を見つけた後で,どの関数との差も平方の積分として表したことが,最小値の証明になっている.

\paragraph{13.15}
微分すると
\[
 x'(t)=1-\cos2t=2\sin^2t,\qquad
 y'(t)=2\sin t\cos t.
\]
$0<t<\pi$では$x'(t)>0$なので,\,$x$は$0$から$\pi$まで狭義に増加する.また$y\geqq0$であり,両端で$y=0$となる.このため曲線は各$x$に対して一つの高さを持ち,曲線の下の領域を$x$方向に積分して重複することはない.

面積$S$は
\[
 S=\int_0^\pi y(t)x'(t)\dd t
 =2\int_0^\pi\sin^4t\dd t.
\]
反転$t\mapsto\pi-t$で両半分の積分は等しい.ウォリスの公式より
\[
 \int_0^{\frac\pi2}\sin^4t\dd t
 =\frac34\frac12\frac\pi2=\frac{3\pi}{16}
\]
なので,
\[
 \boxed{S=4\cdot\frac{3\pi}{16}=\frac{3\pi}{4}}.
\]

曲線の長さ$L$では,根号の中を先にまとめる.
\[
 \begin{aligned}
 \sqrt{x'(t)^2+y'(t)^2}
 &=\sqrt{4\sin^4t+4\sin^2t\cos^2t}\\
 &=2|\sin t|=2\sin t\qquad(0\leqq t\leqq\pi).
 \end{aligned}
\]
従って
\[
 \boxed{L=\int_0^\pi2\sin t\dd t=4}.
\]

回転体の断面は半径$y(t)$の円板であり,\,$x$の増分は$x'(t)\dd t$である.従って体積は
\[
 \begin{aligned}
 V&=\pi\int_0^\pi y(t)^2x'(t)\dd t
 =2\pi\int_0^\pi\sin^6t\dd t\\
 &=4\pi\int_0^{\frac\pi2}\sin^6t\dd t
 =4\pi\cdot\frac{5\pi}{32}
 =\boxed{\frac{5\pi^2}{8}}.
 \end{aligned}
\]
ここでも$x(t)$の単調性により,同じ$x$での円板を二回数えていない.

\paragraph{13.16}
(1) $g(u)=\sin u$の$0$を中心とする六次テイラー多項式は
\[
 P_6(u)=u-\frac{u^3}{6}+\frac{u^5}{120}
\]
である.偶数次の係数は$0$だが,六次までの有限テイラー公式として剰余を評価する.
七階微分は$g^{(7)}(u)=-\cos u$なので,\,$0\leqq u\leqq1$では
$-1\leqq g^{(7)}(u)\leqq0$である.従って
\[
 \sin u=P_6(u)+R_6(u),\qquad
 -\frac{u^7}{5040}\leqq R_6(u)\leqq0.
\]
$u=x^2$を代入して積分した多項式の値は
\[
 Q=\int_0^1\left(x^2-\frac{x^6}{6}+\frac{x^{10}}{120}\right)\dd x
 =\frac13-\frac1{42}+\frac1{1320}
 =\frac{2867}{9240}.
\]
残りの積分は
\[
 -\frac1{5040}\int_0^1x^{14}\dd x
 \leqq A-Q\leqq0
\]
だから
\[
 \boxed{\frac{2867}{9240}-\frac1{75600}
 \leqq A\leqq\frac{2867}{9240}}.
\]
両端は有理数であり,幅$\frac1{75600}$は$\frac1{10000}$より小さい.有限の式に対する区間全体の誤差を積分したので,無限級数と積分の交換は用いていない.

(2) 有限の$R>1$に対して$u=x^2$と置けば
\[
 \int_1^R\sin(x^2)\dd x
 =\frac12\int_1^{R^2}\frac{\sin u}{\sqrt u}\dd u.
\]
右辺で部分積分を行う.\,$T>1$として
\[
 \begin{aligned}
 \int_1^T\frac{\sin u}{\sqrt u}\dd u
 &=\left[-\frac{\cos u}{\sqrt u}\right]_1^T
   -\frac12\int_1^T\cfrac{\cos u}{u^{\cfrac32}}\dd u\\
 &=\cos1-\frac{\cos T}{\sqrt T}
   -\frac12\int_1^T\cfrac{\cos u}{u^{\cfrac32}}\dd u.
 \end{aligned}
\]
境界項$\frac{\cos T}{\sqrt T}$は$0$へ近づく.また
$\cfrac{|\cos u|}{u^{\cfrac32}}\leqq u^{-\frac32}$であり,右辺のべき関数は無限大で積分可能だから,最後の積分は絶対収束する.従って元の広義積分は収束する.

絶対値を付けた場合も同じ置換が使え,
\[
 \int_1^R|\sin(x^2)|\dd x
 =\frac12\int_1^{R^2}\frac{|\sin u|}{\sqrt u}\dd u.
\]
整数$k\geqq1$について,\,$u\in[k\pi,(k+1)\pi]$なら
$\frac1{\sqrt u}\geqq\frac1{\sqrt{(k+1)\pi}}$なので
\[
 \int_{k\pi}^{(k+1)\pi}\frac{|\sin u|}{\sqrt u}\dd u
 \geqq\frac1{\sqrt{(k+1)\pi}}
       \int_{k\pi}^{(k+1)\pi}|\sin u|\dd u
 =\frac2{\sqrt{(k+1)\pi}}.
\]
$R^2\geqq(N+1)\pi$なら区間$[\pi,(N+1)\pi]$は$[1,R^2]$に含まれるので,\,$k=1,\ldots,N$のこの下限を足したものは$\int_1^{R^2}\frac{|\sin u|}{\sqrt u}\dd u$以下である.この下限の和は発散する.実際,\,$t^{-\frac12}$の減少性から
\[
 \sum_{k=1}^N\frac1{\sqrt{k+1}}
 \geqq\int_2^{N+2}\frac1{\sqrt t}\dd t
 =2(\sqrt{N+2}-\sqrt2)\longrightarrow\infty.
\]
従って絶対値の積分は発散し,\,$\int_1^{\infty}\sin(x^2)\dd x$は条件収束する.

\paragraph{13.17}
$n=0$では$C_0=1$である.\,$n\geqq1$で,\,$(1-x^2)^n$を微分し,\,$1$を積分する部分積分を行う.
\[
 \begin{aligned}
 C_n
 &=\left[x(1-x^2)^n\right]_0^1
   +2n\int_0^1x^2(1-x^2)^{n-1}\dd x\\
 &=2n\int_0^1\{1-(1-x^2)\}(1-x^2)^{n-1}\dd x\\
 &=2n(C_{n-1}-C_n).
 \end{aligned}
\]
従って漸化式と有限積の表示は
\[
 \boxed{C_n=\frac{2n}{2n+1}C_{n-1},\qquad
 C_n=\prod_{k=1}^n\frac{2k}{2k+1}}.
\]

極限を求めるため,\,$x=\sin t$と置くと
\[
 C_n=\int_0^{\frac\pi2}\cos^{2n+1}t\dd t.
\]
$W_m=\int_0^{\frac\pi2}\cos^mt\dd t$と書けば$C_n=W_{2n+1}$である.
$t\mapsto\frac\pi2-t$により$W_m$は$\sin$のべきのウォリス積分と等しい.
従って偶数,奇数の積の公式から
\[
 W_{2n}W_{2n+1}=\frac{\pi}{2(2n+1)}.
\]
また$0\leqq\cos t\leqq1$なので$W_{2n+2}\leqq W_{2n+1}\leqq W_{2n}$であり,漸化式を使うと
\[
 \frac{2n+1}{2n+2}
 =\frac{W_{2n+2}}{W_{2n}}
 \leqq\frac{W_{2n+1}}{W_{2n}}\leqq1.
\]
従って隣り合う二つの比は$1$へ近づく.積と比を合わせて
\[
 nC_n^2
 =nW_{2n}W_{2n+1}\frac{W_{2n+1}}{W_{2n}}
 \longrightarrow\frac\pi4.
\]
$C_n>0$なので平方根を取り,
\[
 \boxed{\lim_{n\to\infty}\sqrt n\,C_n=\frac{\sqrt\pi}{2}}.
\]
最後に漸化式から
\[
 C_n-C_{n+1}
 =\left(1-\frac{2n+2}{2n+3}\right)C_n
 =\frac{C_n}{2n+3}.
\]
元の数列の極限をこの式へ戻せば
\[
 \boxed{\lim_{n\to\infty}n^{\frac32}(C_n-C_{n+1})
 =\lim_{n\to\infty}\frac{n}{2n+3}\sqrt n\,C_n
 =\frac{\sqrt\pi}{4}}.
\]

\paragraph{13.18}
まず両端の収束を調べる.\,$-\pi<x<\pi$では$\cos\frac x2>0$だから,半角公式により
\[
 \begin{aligned}
 h_n(x):=\frac{\sin^{2n}x}{\sqrt{1+\cos x}}
 &=\frac{2^{2n}}{\sqrt2}
   \sin^{2n}\frac x2\,\cos^{2n-1}\frac x2.
 \end{aligned}
\]
$n\geqq1$より$2n-1\geqq1$である.\,$x\to\pi-$,\,$x\to-\pi+$のそれぞれで$\cos\frac x2\to0$となり,\,$\sin^{2n}\frac x2$は有界なので,\,$h_n(x)\to0$となる.
さらに$\frac1{1+\e^x}$は両端を含めて連続である.従って元の被積分関数は,各端点で値$0$を補えば連続になる.左右の広義積分はそれぞれ有限に収束し,ここからは対称性や置換を使ってよい.

$h_n$は偶関数である.\,$w(x)=\frac1{1+\e^x}$と置けば$w(x)+w(-x)=1$だから,\,$x=-t$による反転で得た式を元の積分と足して
\[
 \begin{aligned}
 2J_n
 &=\int_{-\pi}^{\pi}h_n(x)\{w(x)+w(-x)\}\dd x\\
 &=\int_{-\pi}^{\pi}h_n(x)\dd x
 =2\int_0^\pi h_n(x)\dd x.
 \end{aligned}
\]
従って指数関数を含む重みは消え,\,$J_n=\int_0^\pi h_n(x)\dd x$となる.

先ほどの半角公式を使い,\,$t=\frac x2$と置くと
\[
 J_n=2^{2n}\sqrt2\int_0^{\frac\pi2}
          \sin^{2n}t\cos^{2n-1}t\dd t.
\]
次に$u=\sin t$と置く.\,$\dd u=\cos t\dd t$なので,\,$\cos$の残りの偶数べきは$(1-u^2)^{n-1}$になり,
\[
 J_n=2^{2n}\sqrt2\int_0^1u^{2n}(1-u^2)^{n-1}\dd u.
\]
最後の積分は多項式の積分である.しかし,展開すると項数が$n$個になるため,ここでは$(1-u^2)$の指数を部分積分で一つずつ下げる.

非負整数$r,s$に対して
\[
 L_{r,s}=\int_0^1u^{2r}(1-u^2)^s\dd u
\]
と置く.\,$s\geqq1$なら
\[
 \begin{aligned}
 L_{r,s}
 &=\left[\frac{u^{2r+1}}{2r+1}(1-u^2)^s\right]_0^1
   +\frac{2s}{2r+1}\int_0^1u^{2r+2}(1-u^2)^{s-1}\dd u\\
 &=\frac{2s}{2r+1}L_{r+1,s-1}.
 \end{aligned}
\]
端点の項は,\,$u=0$では$u^{2r+1}=0$,\,$u=1$では$(1-u^2)^s=0$なので消える.\,$s=0$の初期値は別に
\[
 L_{r,0}=\int_0^1u^{2r}\dd u=\frac1{2r+1}
\]
である.漸化式を$s$回使い,最後にこの初期値へ戻すと
\[
 L_{r,s}
 =\frac{2^s s!}{(2r+1)(2r+3)\cdots(2r+2s+1)}.
\]
$r=n$,\,$s=n-1$を代入し,前の係数$2^{2n}\sqrt2$を掛けることで
\[
 \boxed{J_n=
 \frac{2^{3n-1}\sqrt2\,(n-1)!}
 {(2n+1)(2n+3)\cdots(4n-1)}}.
\]
分母には$n$個の奇数因子がある.例えば$n=1$では分母は$3$だけになり,\,$J_1=\frac{4\sqrt2}{3}$となる.
端点の収束を確かめた後,対称性で指数関数を除き,半角公式と置換で整数べきの積分へ直した.最後の有限積は,その多項式の積分を部分積分で順に結んだ結果である.
\paragraph{13.19}
(1) $f(x)=x^2-2$に対して,点$(x_n,f(x_n))$における接線と$x$軸の交点の$x$座標を求めると,
\[
 x_n-\frac{f(x_n)}{f'(x_n)}
 =x_n-\frac{x_n^2-2}{2x_n}
 =\frac12\left(x_n+\frac2{x_n}\right)
\]
である.また$x_1>\sqrt2$であり,\,$x_{n+1}-\sqrt2=\frac{(x_n-\sqrt2)^2}{2x_n}>0$だから,すべての項は$\sqrt2$より大きく,分母は$0$にならない.

(2) $r=\sqrt2$と略記する.分子・分母をそれぞれ平方にすると,
\[
 b_{n+1}
 =\frac{x_n^2+2-2rx_n}{x_n^2+2+2rx_n}
 =\left(\frac{x_n-r}{x_n+r}\right)^2=b_n^2.
\]
$b_1=q$なので,順に$b_2=q^2,b_3=q^4,\ldots$となる.よって
\[
 b_n=q^{2^{n-1}},\qquad
 x_n=r\frac{1+b_n}{1-b_n}
 =\sqrt2\frac{1+q^{2^{n-1}}}{1-q^{2^{n-1}}}.
\]
(3) $0<q<1$より$b_n\to0$であり,
\[
 \frac{x_n-\sqrt2}{q^{2^{n-1}}}
 =\frac{2\sqrt2}{1-q^{2^{n-1}}}\longrightarrow2\sqrt2.
\]
誤差が小さくなる様子を,漸化式の大小関係だけでなく一般項からも確かめられた.

\paragraph{13.20}
(1) 接線は$y=2ax-a^2$,\,$y=2bx-b^2$だから,交点は
\[
 T\left(\frac{a+b}{2},ab\right)
\]
である.弦$AB$の傾きは$\frac{b^2-a^2}{b-a}=a+b$である.点$(x,x^2)$での接線の傾きは$2x$なので,傾きが$a+b$となる接点は$x=\frac{a+b}{2}$の点,すなわち$C$に限る.

(2) $h=b-a$とおく.弦の式は$y=(a+b)x-ab$なので,
\begin{align*}
 Q&=\int_a^b(x-a)(b-x)\dd x\\
 &=\int_0^h u(h-u)\dd u=\frac{h^3}{6}.
\end{align*}
ここで$u=x-a$と置いた.弦の中点を$D$とすると,\,$D,C,T$はいずれも$x=\frac{a+b}{2}$上にある.縦の長さは
\[
 DC=\frac{a^2+b^2}{2}-\frac{(a+b)^2}{4}=\frac{h^2}{4},
 \qquad DT=\frac{a^2+b^2}{2}-ab=\frac{h^2}{2}.
\]
三角形$ABC$を$CD$で分け,\,$CD$を共通の底辺として左右の面積を足せば,\,$P=\frac12\cdot DC\cdot h=\frac{h^3}{8}$となる.同様に$R=\frac{h^3}{4}$である.したがって$Q:P:R=4:3:6$となる.取り尽くし法で得た$Q=\frac43P$を,今度は積分で確かめられた.

(3) $T$の$x$座標は$c$で一定である.\,$a=c-\frac h2,b=c+\frac h2$なので,\,$T$の$y$座標は$c^2-\frac{h^2}{4}$となる.\,$h>0$を自由に取れるから,軌跡は直線$x=c$上の$y<c^2$の部分である.端の点$(c,c^2)$は含まない.

\paragraph{13.21}
(1) 二項定理を使うと,
\[
 f_n(x)=\frac1{n!}\sum_{j=0}^n(-1)^j\binom nj x^{n+j}.
\]
$k$回微分して$x=0$を代入すると,\,$k<n$なら$0$,\,$k=n+j$なら
\[
 f_n^{(k)}(0)=(-1)^j\binom nj\frac{(n+j)!}{n!}
\]
となり,整数である.\,$f_n(1-x)=f_n(x)$を$k$回微分すれば,\,$f_n^{(k)}(1)=(-1)^k f_n^{(k)}(0)$なので,こちらも整数である.

(2) 部分積分を繰り返すと,
\[
 I_n=\left[\e^x\sum_{k=0}^{2n}(-1)^k f_n^{(k)}(x)\right]_0^1.
\]
各回で$f_n$を微分し,\,$\e^x$を積分した.次数が$2n$なので,\,$2n+1$回微分すると$0$になり,残る積分はなくなる.(1)より両端の和は整数であるから,\,$I_n=A_n\e+B_n$と表せる.

(3) $0<x<1$で被積分関数は正なので$I_n>0$である.また$x(1-x)\le\frac14$,\,$\e^x\le\e$より
\[
 I_n\le\int_0^1\frac{\e}{4^n n!}\dd x=\frac{\e}{4^n n!}\longrightarrow0.
\]
もし$\e=\frac ab$と書けるとする.ただし$a,b$は正の整数である.すると
\[
 bI_n=A_na+bB_n
\]
は整数であり正である.ところが$I_n\to0$だから,十分大きな$n$では$0<bI_n<1$となる.この範囲に整数は存在しないので矛盾し,\,$\e$は無理数である.

ここで積分の値を具体的に求める必要はなかった.部分積分は値の表し方を,大小比較は値の大きさを与え,二つを合わせることで整数についての矛盾が得られた.

\paragraph{13.22}
(1) $\frac1{m+k}=\int_0^1x^{m+k-1}\dd x$だから,有限個の和を積分の中へ入れて
\begin{align*}
 S_{m,n}
 &=\int_0^1x^{m-1}\sum_{k=0}^n(-1)^k\binom nk x^k\dd x\\
 &=\int_0^1x^{m-1}(1-x)^n\dd x
 =\frac{(m-1)!n!}{(m+n)!}.
\end{align*}
最後の値は整数指数のベータ積分である.公式を使わずに求めるなら,\,$n\ge1$で部分積分して
\[
 \int_0^1x^{m-1}(1-x)^n\dd x
 =\frac nm\int_0^1x^m(1-x)^{n-1}\dd x
\]
を$n$回用いればよい.最後は$\int_0^1x^{m+n-1}\dd x=\frac1{m+n}$となり,同じ階乗の式を得る.

(2) 階乗の比を積に戻すと,
\[
 n^m S_{m,n}
 =(m-1)!\prod_{j=1}^m\frac{n}{n+j}\longrightarrow(m-1)!.
\]
$m$は固定されているので,積の因子の個数も固定されている.元の和では項数が増え,正負の項が打ち消し合うが,積分でまとめると極限を求めやすい形になった.

\paragraph{13.23}
(1) 直接積分して,
\[
 \int_0^1p(x)\dd x=\frac13-\frac12+\frac16=0,
 \qquad
 \int_0^1xp(x)\dd x=\frac14-\frac13+\frac1{12}=0.
\]
(2) $u=1-a$,\,$v=-\frac16-b$とおくと,\,$x^2-(ax+b)=p(x)+ux+v$である.平方を展開して積分し,(1)を使うと,
\begin{align*}
 E(a,b)&=\int_0^1p(x)^2\dd x+\int_0^1(ux+v)^2\dd x\\
 &=\frac1{180}+\int_0^1(ux+v)^2\dd x\ge\frac1{180}.
\end{align*}
ここで$p(x)^2=x^4-2x^3+\frac43x^2-\frac13x+\frac1{36}$を積分した.等号は連続な非負関数$(ux+v)^2$の積分が$0$のときであり,\,$ux+v=0$が全区間で成り立つ場合に限る.よって$u=v=0$,すなわち$a=1,b=-\frac16$のとき,最小値$\frac1{180}$を取る.

(3) ベクトルの和は成分ごとに足すので,$|\boldsymbol u+\boldsymbol v|^2=(u_1+v_1)^2+(u_2+v_2)^2$を展開すれば,問題文の等式を得る.内積が$0$となる関係を\textbf{直交}という.どちらのベクトルも零ベクトル$(0,0)$でない場合,これは二つの矢印が垂直になることを表す.実際,差$\boldsymbol u-\boldsymbol v=(u_1-v_1,u_2-v_2)$を使うと,$|\boldsymbol u-\boldsymbol v|^2=|\boldsymbol u|^2+|\boldsymbol v|^2$となるので,三平方の定理の逆を使えばよい.

関数$f,g$に対する$\int_0^1f(x)g(x)\dd x$が,ベクトルの内積に対応している.(1)は$p$が$1$と$x$のそれぞれに直交することに当たり,したがって任意の一次式$ux+v$に対しても積の積分が$0$になる.その結果,平方を展開した際の交差する項が消えた.ベクトルについての「直交していれば長さの二乗が足される」という計算を,関数にも同じ形で使えたのである.

\paragraph{13.24}
(1) 確率は正方形と円板$x^2+y^2\le r^2$が重なる部分の面積である.\,$0\le r\le1$なら四分円なので$\frac{\pi r^2}{4}$である.\,$1<r\le\sqrt2$では$a=\sqrt{r^2-1}$とおくと,\,$0\le x\le a$で正方形の高さ全体が入り,\,$a<x\le1$では高さが$\sqrt{r^2-x^2}$となる.よって
\[
 F(r)=a+\int_a^1\sqrt{r^2-x^2}\dd x.
\]
$x=r\sin\theta$で積分すれば,原始関数は
\[
 \frac12\left\{x\sqrt{r^2-x^2}+r^2\arcsin\frac xr\right\}.
\]
ここで$\arcsin$の値域は$\left[-\frac\pi2,\frac\pi2\right]$である.\,$\arcsin\frac ar=\frac\pi2-\arcsin\frac1r$と$\sqrt{r^2-a^2}=1$より,
\[
 F(r)=\sqrt{r^2-1}+r^2\left(\arcsin\frac1r-\frac\pi4\right)
 \quad(1<r\le\sqrt2).
\]
$r=1$で$\frac\pi4$,\,$r=\sqrt2$で$1$となり,隣の範囲とも一致する.

(2) 被積分関数は$x,y$を交換しても変わらず,二つの三角形もその交換で移り合う.したがって正方形全体の平均は,下側の三角形$0\le y\le x\le1$での積分の$2$倍になる.この積分の意味は,小さい長方形ごとに「距離$\times$面積」を足していくことである.正方形の面積が$1$なので,さらに全体の面積で割る必要はない.

内側の積分では$x>0$を固定し,\,$y=xt$と置く.すると
\[
 \int_0^x\sqrt{x^2+y^2}\dd y
 =x^2\int_0^1\sqrt{1+t^2}\dd t.
\]
$x=0$では両辺とも$0$である.右辺の$t$についての積分は$x$によらないから,
\begin{align*}
 M&=2\int_0^1x^2\dd x\int_0^1\sqrt{1+t^2}\dd t\\
 &=\frac23\cdot\frac12\{\sqrt2+\log(1+\sqrt2)\}\\
 &=\frac{\sqrt2+\log(1+\sqrt2)}3.
\end{align*}
距離の平方の平均なら$\int_0^1\int_0^1(x^2+y^2)\dd y\dd x=\frac23$だが,その平方根が距離の平均になるわけではない.平均する量自体を積分することが必要である.
